% ── Part I, Chapter 1 ────────────────────────────────────────
% The Interior Field
% Jeanne-Marie Guyon, France, 1688

\chapter{The Room That Would Not Stop}

\strandmarker{Guyon}{France, 1688}

\lettrine{T}{here} is a window in the room that gives onto a garden, but the window
is not what she looks at. She has come to understand this about
herself, which is more than most people manage.

The window is a fact. The garden beyond it---the clipped yew, the
gravel paths, the fountain that someone has let run---is also a
fact, gathered and arranged and held in place by a logic she did not
choose. What passes through her is not the garden. It is something
that the garden might, on a different day, nearly resemble.

She has written many things today. She writes the way water finds a
level: not by intention but by the exhaustion of all other
possibilities. The pages accumulate on the desk to her left. She
does not reread them. Rereading would be a form of interference, a
hand reaching back into a process that has already moved past the
point where hands are useful.

Fénelon will ask, when she sees him next, whether she has been
sleeping. She will tell him that sleep is a different country with
different customs and she has not been issued the correct
documents.

He will not laugh at this. He will make the face he makes when
something that should trouble him has decided instead to interest
him: a small displacement around the eyes, a settling of the jaw.
She finds this face useful in the way that a reliable instrument
is useful. She trusts its readings precisely because she has seen
it fail to adjust for her.

\scenebreak

Outside the window the fountain runs. She can hear it when the house
is quiet enough, which is not often. There are people in the next
room conducting an argument about property she will never understand
and does not attempt to. Property is what happens when attachment
meets administration. She has been trying to loosen her grip on
both.

This is not a spiritual exercise, whatever anyone says. It is more
like what a hand does when it has been holding something too tightly
for too long and finally begins to understand that the grip itself
is the problem.

The pages. The pen still wet. The window full of garden and the
garden full of particular, impermanent light.

She writes: \textit{The soul that has ceased to insist on its own motion
finds that it has not ceased to move.}

She looks at this for a long time without reading it.

Then she writes beneath it: \textit{This is either the most important
thing I have said or the most useless. I cannot tell which and I
suspect the question is misformed.}

\scenebreak

The argument in the next room has concluded or migrated. Either way
the house has quieted. The fountain is more present now, its
particular rhythm---not quite regular, altered by wind and the
angle of supply---more audible.

She will not tell Bossuet about the pages. This is not deception.
It is a different kind of accuracy: she has learned that certain
processes, examined too early by the wrong instrument, will
collapse back into their components and lose whatever coherence
they were developing. Bossuet is not the wrong instrument for
everything. He is the wrong instrument for this.

She thinks: there are forms of knowledge that require the knower
to get out of the way.

She thinks: I am still in the way.

She picks up the pen.

